
Automated health metrics detect dataset shifts, context-aware graphs recommend task-specific successors, and load-time instrumentation measures adoption. Proof-of-concept validation demonstrates that all three components work at mechanistic scale: 88.3\% precision in deprecation detection, 100\% accuracy on synthetic task inference, and <5\% overhead for adoption tracking. Deployment requires moving from proof-of-concept validation to production-scale efficacy testing.

Formal dataset deprecation mechanisms demonstrate mechanistic feasibility, enabling deployment-scale testing. The three-component system fills gaps that software package managers (lack task-aware succession) and dataset documentation (lack executable enforcement) cannot solve in isolation. Sub-hypothesis decomposition validated each mechanism independently before integration, establishing mechanistic feasibility for HuggingFace-compatible infrastructure.

\subsubsection{7.1 Future Work}

\textbf{Immediate:} Phase 5 baseline comparison (measure adoption lift versus informal mechanisms through controlled experiment), real HuggingFace API integration (validate generalization beyond mock data), cross-platform deployment (OpenML, Kaggle, Papers with Code).

\textbf{Long-term:} Community-maintained successor graphs (crowdsourced task-specific recommendations), automated deprecation workflows triggered by health metrics (reduce maintainer burden), integration with model deployment pipelines (upstream-downstream dependency tracking), longitudinal efficacy studies (A/B testing deprecation interventions).

The path from mechanistic feasibility to measured efficacy is clear. Validated infrastructure enables efficacy measurements that current practice cannot perform: quantitative adoption lift, intervention optimization, and long-term reproducibility impact. Formal deprecation mechanisms are demonstrated feasible; deployment-scale testing remains.
